\honest{Making Beliefs Pay Rent (in Anticipated Experiences)}{Eliezer Yudkowsky}{2007}

I open with the old riddle: if a tree falls in a forest and no one hears it, does it make a sound? One person says yes, because it makes vibrations in the air. The other says no, because no one hears it.

The foundational skill of rationality, I say, ``might be'' the ability to spot in your own head whether you have a map of something or not. I do not yet say how.

Would the two arguers expect to see the tree fallen in different places, or hear anything different on a recorder? No. So, I say, their maps of the world do not differ. They do disagree about something: what the word ``sound'' means. I treat an argument about a word as the same failure as a belief with no content, until near the end.

It would be tempting, I say, to allow only beliefs about direct sensations. But atoms and floors are real, though we never sense them directly. This is right.

Then comes the best example in the essay. You drop a bowling ball from the top of a building and ask on which tick of the clock you will hear it land. To answer, you need two beliefs: that gravity is 9.8 meters per second per second, and that the building is about 120 meters tall. Neither predicts anything alone. Together they predict a crash five seconds later. So beliefs make predictions in combination. I will forget this by the end.

Humans can model what they cannot see, I say, and can also believe in what is not there. My example is that alchemists believed phlogiston caused fire, and the belief predicted nothing in advance. Phlogiston was a chemists' theory, and it did make predictions: some metals gained weight when burned, and that is how the theory was refuted.

Next, an English professor tells you that a writer called Wulky Wilkinsen is a ``retropositional author,'' because his books show ``alienated resublimation,'' and each term is defined only by the other. What should you expect from his books? ``Nothing.'' I made up both terms. The beliefs I trust in this essay come from physics. The one empty belief I invent comes from a literature class.

I call beliefs that are connected only to each other ``floating'' beliefs, and I say, with no evidence, that they are ``a uniquely human flaw among animal species.''

Empiricism, I say, means constantly asking what experiences a belief predicts or, better, forbids. This is the actual test, and it arrives three-quarters of the way through. My example of a belief that forbids nothing is élan vital, a life force that nobody believes in any more. I do not apply the test to any belief that anyone holds today.

If you cannot find a difference in what the two sides expect, I say, you are arguing either about labels or about floating beliefs. These are different problems with different cures, and I give one cure.

I end: ``Every question of belief should flow from a question of anticipation.'' Not every claim in mathematics, ethics or history predicts a sensation, and the bowling ball showed that beliefs predict in groups. ``If a belief turns deadbeat, evict it.''
